\section{Trapped domains, collars, and weak supports}
\label{n3:domains:section}

We prepare the domain and boundary values for the elliptic deformation.
The construction has three steps: choose maximal trapped regions,
foliate their outer collars, and fix small expansion offsets.
Throughout, the data satisfy
Proposition~\ref{n3:reduction:rest-reduction}.  Write $S_0$ for their inner
boundary.  Thus $S_0$ is strictly future outer trapped, $K$ has compact
support, and the single end has the smooth $O(r^{-1})$ bounds obtained in
that reduction.  In particular, sufficiently large coordinate spheres
have positive outward mean curvature.  The function
\begin{equation}\label{n3:domains:margin}
 \mathfrak m(x)=8\pi\bigl(\mu(x)-|J(x)|_g\bigr)
\end{equation}
is positive, has a positive lower bound on every compact set, and has
the positive $r^{-3-\delta}$ lower bound on the end supplied by the
reduction, with $0<\delta<1$.  Let $A_*$ denote the reduced data's
filled enclosing perimeter infimum.

\subsection{The bounding-region convention}

For a symmetric tensor $Q$ and a two-sided surface $\Sigma$ with a
specified unit normal $\nu$, put
\begin{equation}\label{n3:domains:expansion}
 \Theta_Q(\Sigma,\nu)
   =H_\Sigma(\nu)+\tr_\Sigma Q.
\end{equation}
For the boundary of a domain the specified normal always points
\emph{out of that domain}.  Notice both elementary identities
\begin{equation}\label{n3:domains:shift-reversal}
 \Theta_{Q-(c/2)g}=\Theta_Q-c,
 \qquad
 \Theta_{-Q}(\Sigma,-\nu)=-\Theta_Q(\Sigma,\nu).
\end{equation}
The factor $1/2$ in the first formula is the reciprocal of the surface
dimension.

We use the following geometric theorem.  Its boundary parts may be
disconnected.

\begin{theorem}[MOTS barriers and the total trapped region]
\label{n3:domains:region-theorem}
Let $(N,g,Q)$ be a smooth compact three-dimensional initial data set
with boundary $\Gamma_-\sqcup\Gamma_+$.  Assume $\Gamma_+$ is nonempty
and has positive expansion with its normal out of $N$.  The inner part
$\Gamma_-$ may be empty; when it is present, assume its expansion is
negative with its normal into $N$.

Consider all smooth domains whose inner boundary is $\Gamma_-$ and
whose remaining boundary $\Sigma$ lies in the interior of $N$ and
satisfies $\Theta_Q(\Sigma)\leq0$.  They avoid $\Gamma_+$, and their
normal on $\Sigma$ points out of the domain.  If their union is
nonempty, its closure has a smooth compact embedded stable MOTS
frontier in the interior, separating it from $\Gamma_+$.  This closure
is itself the closure of such a domain and contains every domain in
the defining class.  If $\Gamma_-$ is nonempty, the two strict barriers
also give a nonempty smooth embedded enclosing MOTS.
\end{theorem}

The total-region assertion is
\cite[Definitions 7.1--7.2 and Theorem 7.3]{AnderssonMetzger2009}; the
barrier assertion is Theorem 3.1 there, with stability furnished by
Theorem 4.1.  The preliminary conventions in Section 2 of that paper
allow an empty inner boundary.  These results require no energy
condition on $Q$.  We therefore apply them to the shifted tensors in
Equation~\eqref{n3:domains:shift-reversal}, even though these tensors
need not satisfy the dominant energy condition.  A component of the
complement of a smooth trapped domain which meets no designated outer
boundary may be filled: its entire boundary is then deleted and the
remaining boundary has the same expansion.  The maximal region is
therefore filled in this relative sense.

We use compact inner fillings only to express enclosing competitors.
One may attach a smooth collar behind $S_0$, extend the one-sided
metric smoothly, and complete the inner side by a fixed compact
filling.  All competitors contain that filling and only their
boundaries in the original exterior are counted.  This introduces no
geometric condition behind $S_0$.  For perimeter minimization with
this smooth obstacle, the relevant regularity input is
\cite[Regularity Theorem 1.3(iii)]{HuiskenIlmanen2001}: the minimizing
boundary is $C^{1,1}$ and is smooth away from contact.  We retain the
filled enclosing class and its smooth approximation from
Section~\ref{n3:reduction:section}.  The trapped regions below are not
assumed to minimize perimeter.

\subsection{Choosing the two maximal regions}

Fix a large sphere $S_{R_0}$ outside $\supp K$ such that all coordinate
spheres with radius at least $R_0$ have positive outward mean
curvature.  A compact smooth domain with $\Theta_{\pm K}\leq c\leq0$
cannot reach this part of the end.  Indeed, at a point where its
boundary maximizes the coordinate radius, it lies inside the tangent
coordinate sphere and has the same outward normal.  The graph
second-derivative comparison gives
$H_{\partial D}\geq H_{S_r}>0$ there; also $K=0$ there.  This
contradicts its assumed expansion.  The same argument applies to an
enclosing domain with prescribed inner obstacle, since the maximum
radius is attained on its free boundary.  Consequently all the
negative-threshold regions used below lie in one fixed compact radial
range.  They may be constructed using any sufficiently distant
sphere as the outer barrier, without dependence on that choice.

For
\begin{equation}\label{n3:domains:black-interval}
 \max_{S_0}\Theta_K(S_0)<c<0,
\end{equation}
let $\mathcal B_c$ be the closed total trapped region for the
condition $\Theta_K\leq c$, with $S_0$ as the required inner
boundary and with its inner filling understood.  It is nonempty:
a sufficiently short outward collar of $S_0$ supplies a strict seed.
Theorem~\ref{n3:domains:region-theorem}, applied to $K-(c/2)g$, gives
its smooth stable frontier with expansion exactly $c$.  The regions
$\mathcal B_c$ increase with $c$.  Their complement toward the end
has no bounded component, since such a component could be filled to
enlarge the total trapped region.

We will choose $c_b<0$ close to zero and set
$\mathcal B=\mathcal B_{c_b}$.  In the complement of $\mathcal B$,
use the full black frontier and a distant coordinate sphere as
\emph{outer} barrier faces.  With the normal out of this complement,
the black frontier has expansion for $-K$ equal to $-c_b>0$.
For $c<0$ close to zero, let $\mathcal W_c$ be the closed total
trapped region of compact smooth domains in that complement with
\begin{equation}\label{n3:domains:white-condition}
 \Theta_{-K}=H-\tr_\Sigma K\leq c.
\end{equation}
There is no required inner boundary in this application of
Theorem~\ref{n3:domains:region-theorem}.  The shifted expansion of a
reversed black face is $-c_b-c>0$, and that of a distant sphere is
also positive.  Thus the theorem applies whenever the defining union
is nonempty.  When the union is empty we simply put
$\mathcal W_c=\varnothing$.  When nonempty, it has smooth stable
frontier of white expansion $c$ and lies a positive distance from all
black faces.  The family $\mathcal W_c$ is increasing, with the black
region held fixed.

To control limits of the deformation, we choose thresholds at which these
regions are right-continuous.  The following lemma includes empty regions.

\begin{lemma}[Right-continuous thresholds]\label{n3:domains:right-continuity}
Let $E_c$, $c$ in an interval, be an increasing family of closed
subsets of a fixed compact metric space.  Outside a countable set of
parameters,
\begin{equation}\label{n3:domains:hausdorff-right}
 E_{c_j}\longrightarrow E_c\quad\hbox{in Hausdorff distance whenever}
 \quad c_j\downarrow c.
\end{equation}
Here continuity at $E_c=\varnothing$ means that $E_{c'}$ is empty
for every sufficiently close $c'>c$.
\end{lemma}

\begin{proof}
Choose a countable dense set $\{x_i\}$ and a number $L$ larger than
the diameter of the compact space.  Set
$d_c(x)=\min\{L,\operatorname{dist}(x,E_c)\}$, with the value $L$
when $E_c$ is empty.  For each $i$, $c\mapsto d_c(x_i)$ is a
bounded nonincreasing real function, so its right discontinuities
form a countable set.  Avoid their countable union.  Then
$d_{c_j}(x_i)\to d_c(x_i)$ for all $i$ whenever $c_j\downarrow c$.
Every $d_c$ is $1$-Lipschitz.  A finite net chosen from the dense set
therefore upgrades this convergence to uniform convergence on the
compact space.

For nonempty $E_c$, if points $y_j\in E_{c_j}$ stayed a fixed
distance from $E_c$, uniform convergence evaluated at $y_j$ would
contradict $d_{c_j}(y_j)=0$.  The opposite Hausdorff inclusion follows
from $E_c\subset E_{c_j}$.  For empty $E_c$, uniform convergence to
$L$ is incompatible with any nonempty $E_{c_j}$, since its distance
function has a zero.  This proves the stated convention as well.
\end{proof}

Choose $c_b$ at such a right-continuity point of the black family,
then choose $c_w<0$ at such a point of the white family for this
fixed black region.  Both choices can be made arbitrarily close to
zero.  Write $\mathcal W=\mathcal W_{c_w}$, and let $\Omega$ be the
closure of the component of the complement of
$\mathcal B\cup\mathcal W$ which reaches the selected end.  Its
boundary $B$ is a finite disjoint union of complete components of
the two smooth frontiers.  Write
\begin{equation}\label{n3:domains:faces}
 B=B_b\sqcup B_w,\qquad
 H=H_B(\nu),\qquad P_B=\tr_BK,
\end{equation}
where $\nu$ points into $\Omega$.  Thus
\begin{equation}\label{n3:domains:face-expansions}
 H+P_B=c_b\quad\hbox{on }B_b,
 \qquad
 H-P_B=c_w\quad\hbox{on }B_w.
\end{equation}
The set $B_w$ may be empty, and $B_b$ need not contain every black
frontier component: a white component may cut a black face off from
the end.  We retain exactly the boundary of the component that reaches
infinity.  The compact, disjoint frontiers create no intersection corners.

The excluded side of $B$ contains $S_0$ and a collar of it.  Hence
every compact smooth enclosing cut $T$ in $\Omega$, with all its
components retained, is an enclosing competitor for the original
filled obstacle after adjoining the discarded regions to its inner
side.  It follows that
\begin{equation}\label{n3:domains:inherited-area}
 |T|_g\geq A_*.
\end{equation}
The same conclusion holds for perimeter competitors.  This is a
statement about inclusion of competitor classes; no minimizing
property of $B$ is used.

\subsection{Outward collars, including zero stability eigenvalue}

\begin{lemma}[Collars of a maximal frontier]\label{n3:domains:stable-collar}
Let $\Sigma$ be a connected component of the frontier of a closed
total trapped region at threshold $c$, for a fixed tensor $Q$.
There is a smooth foliation $\Sigma_s$, $0\leq s<s_0$, of an outward
collar, with $\Sigma_0=\Sigma$ and positive outward normal velocity,
such that
\begin{equation}\label{n3:domains:collar-expansion}
 \Theta_Q(\Sigma_s)>c\quad(0<s<s_0),\qquad
 \Theta_Q(\Sigma_0)=c.
\end{equation}
In particular its leaf parameter $s$ is a smooth defining function
and $|\nabla s|$ is bounded above and below by positive constants.
\end{lemma}

\begin{proof}
Use normal graphs $\Sigma(v)$ over $\Sigma$ and pull their expansion
back to $\Sigma$.  The map
\[
 \Phi:C^{2,\alpha}(\Sigma)\longrightarrow C^{0,\alpha}(\Sigma),
 \qquad
 \Phi(v)=\Theta_Q(\Sigma(v))-c
\]
is smooth near zero and has linearization $L=D\Phi(0)$, the MOTS
stability operator for $Q-(c/2)g$.  Its principal part is
$-\Delta_\Sigma$.  Stability gives a principal eigenvalue
$\lambda\geq0$ and a positive smooth eigenfunction $\varphi$.
If $\lambda>0$, take $v=s\varphi$.  Smooth dependence gives
$\Phi(s\varphi)=s\lambda\varphi+O(s^2)$ in $C^0$, which is
positive for all sufficiently small $s>0$.  The graph velocity is
positive as well.

Suppose $\lambda=0$.  The kernel of $L$ is spanned by $\varphi$,
and the adjoint kernel is spanned by a positive smooth function
$\varphi^*$.  These are the principal-eigenfunction facts for a
scalar elliptic operator on a connected closed surface.  For
example, simplicity of the kernel also follows by writing a kernel
element as $v=\varphi w$ and applying the strong maximum principle
to the resulting equation for $w$, whose zeroth-order term vanishes.
Consider the augmented map
\[
 \mathcal G(v,a)
   =\left(\Phi(v)-a,\ \int_\Sigma v\,\dd A\right).
\]
Its derivative at $(0,0)$ is
\begin{equation}\label{n3:domains:augmented-linearization}
 (v,a)\longmapsto
 \left(Lv-a,\ \int_\Sigma v\,\dd A\right).
\end{equation}
This map is an isomorphism.  Indeed, to solve $Lv-a=f$ with
$\int v=b$, the Fredholm compatibility condition uniquely fixes
\[
 a=-\frac{\int_\Sigma\varphi^*f\,\dd A}
          {\int_\Sigma\varphi^*\,\dd A}.
\]
The equation for $v$ then has a solution; adding a multiple of
$\varphi$ imposes its prescribed integral uniquely.  Elliptic
estimates and this one-dimensional normalization give a bounded
inverse between the indicated H\"older spaces.

The implicit-function theorem consequently supplies smooth
$v(s),a(s)$ with $\mathcal G(v(s),a(s))=(0,s)$.  Differentiating at
zero and pairing with $\varphi^*$ gives $a'(0)=0$ and
\[
 v'(0)=\frac{\varphi}{\int_\Sigma\varphi\,\dd A}>0.
\]
After shortening the interval, all these graphs form an outward
foliation and have constant expansion $c+a(s)$.  If $a(s)\leq0$
at any positive $s$ in this interval, replace only this frontier
component by $\Sigma(v(s))$.  The collar is disjoint from all other
components and designated barriers, so the new smooth bounding
domain strictly contains the maximal region and its every free
boundary component has expansion at most $c$.  This contradicts
maximality.  Therefore $a(s)>0$ for every such $s>0$.

In either case positive graph velocity and compactness give the
asserted bounds on $|\nabla s|$.  Smooth elliptic bootstrapping of
the graph equation gives smooth leaves in the zero-eigenvalue case.
\end{proof}

Apply this lemma separately with $Q=K$ to black faces and $Q=-K$
to white faces.  There are finitely many faces, so choose disjoint
collars in $\Omega$ and shorten them to have a common positive
parameter range when convenient.  Write $\nu_s=\nabla s/|\nabla s|$
for their outward leaf normal; it agrees with $\nu$ at $B$.

\subsection{From weak exterior supports to a smooth enclosure}

The sublevel limits used below need not have smooth frontiers.  We now
show how an expansion bound on their exterior supports produces a smooth
enclosure.  The argument applies to arbitrary closed sets, without
positive reach or a curvature bound on the supporting surfaces.

\begin{definition}\label{n3:domains:support-definition}
Let $D$ be closed in a smooth three-manifold and let $x\in\partial D$.
A smooth exterior support is a smooth function $\phi$ near $x$ with
$\phi(x)=0$, $\dd\phi(x)\ne0$, and
$D\subset\{\phi\leq0\}$ locally near $x$.  Its normal is
$n=\nabla\phi/|\nabla\phi|$.  Define
\begin{equation}\label{n3:domains:test-expansion}
 \mathcal E_Q[\phi](x)
   =\frac{(g^{ij}-n^in^j)\nabla_i\nabla_j\phi}{|\nabla\phi|}
       +\tr Q-Q(n,n).
\end{equation}
We say $D$ has exterior-support expansion at most $c$ if every such
support at every point of its frontier has
$\mathcal E_Q[\phi]\leq c$.
\end{definition}

For a smooth domain this is exactly the usual upper bound on its
outward expansion: at contact an exterior supporting graph has mean
curvature at most that of the enclosed smooth graph.  The definition
is unchanged by increasing smooth reparametrizations of $\phi$,
because the additional Hessian term is a multiple of
$\dd\phi\otimes\dd\phi$ and has zero tangential trace.

\begin{lemma}[Parallel neighborhoods and additive error]
\label{n3:domains:parallel-support}
Suppose the frontier of a closed set $D$ lies in a fixed compact
interior region of smooth $(N,g,Q)$ and has exterior-support
expansion at most $c$.  Put
$D_z=\{y:\operatorname{dist}(y,D)\leq z\}$.  For sufficiently small
$z>0$, all exterior supports of $D_z$ have expansion at most
\begin{equation}\label{n3:domains:parallel-estimate}
 c+Cz.
\end{equation}
The allowable radius and $C$ depend only on the compact ambient
geometry and $Q$, and not on the support, its second fundamental
form, or the regularity of $D$.
\end{lemma}

\begin{proof}
Take a support $\phi$ to $D_z$ at $x'$.  There is a nearest point
$x\in D$ and a length-$z$ minimizing unit-speed geodesic
$\gamma:[0,z]\to N$ from $x$ to $x'$.  For the radii under
consideration the entire segment lies in a fixed smooth compact
neighborhood and is shorter than its injectivity radius.  The ball
$\overline B_z(x)$ is contained in $D_z$ and touches its exterior
support at $x'$.  Its tangent plane there therefore agrees with that
of the support and
\begin{equation}\label{n3:domains:radial-normal}
 \frac{\nabla\phi(x')}{|\nabla\phi(x')|}=\dot\gamma(z).
\end{equation}
Let $P_t:T_xN\to T_{\gamma(t)}N$ denote parallel transport and put
$P=P_z$.

For every $X\in T_xN$ solve the Jacobi boundary problem
\begin{equation}\label{n3:domains:jacobi-endpoints}
 J_X''+\operatorname{Rm}(J_X,\dot\gamma)\dot\gamma=0,
 \qquad J_X(0)=X,\qquad J_X(z)=PX.
\end{equation}
Here primes denote covariant differentiation along $\gamma$.
The short interval and absence of conjugate points make this problem
uniquely solvable.  Its estimates can be seen without division by an
uncontrolled quantity.  In a parallel frame write
$j_X(t)=P_t^{-1}J_X(t)$ and $\mathcal R(t)$ for the curvature
coefficient.  The integral equation is
\[
 j_X(t)=X+tA_X-
       \int_0^t(t-s)\mathcal R(s)j_X(s)\,\dd s,
 \qquad
 A_X=\frac1z\int_0^z(z-s)\mathcal R(s)j_X(s)\,\dd s.
\]
It implies, by absorption for small $z$,
\begin{equation}\label{n3:domains:jacobi-estimates}
 \sup_{0\leq t\leq z}|j_X(t)|\leq2|X|,
 \qquad |A_X|\leq Cz|X|.
\end{equation}
Also $\langle J_X,\dot\gamma\rangle$ is affine in $t$ and has
equal endpoint values.  Thus
$\langle A_X,\dot\gamma(0)\rangle=0$.

We may therefore prescribe the first jet of a smooth unit vector
field $V$ near $x$ by
\[
 V(x)=\dot\gamma(0),\qquad \nabla_XV(x)=A_X.
\]
To realize it with uniform second-jet bounds, choose a smooth local
orthonormal frame obtained by radial parallel transport from $x$.
In that frame take the affine coefficient vector with the prescribed
value and first derivatives, and normalize its length.  The
orthogonality just proved means normalization preserves its prescribed
first derivatives.  The frame has uniformly bounded derivatives on
a fixed small coordinate ball, and Equation~\eqref{n3:domains:jacobi-estimates}
therefore gives uniformly bounded second derivatives of $V$.
This construction uses the radial direction and the ambient
geometry, with no derivatives of the support $\phi$.

Define the local endpoint map
\[
 F_z(y)=\exp_y\bigl(zV(y)\bigr).
\]
Its differential at $x$ is the terminal value of the Jacobi field
with initial values $X,\nabla_XV(x)$.  By
Equation~\eqref{n3:domains:jacobi-endpoints}, it has the exact property
\begin{equation}\label{n3:domains:exact-endpoint-isometry}
 F_z(x)=x',\qquad (\dd F_z)_x=P,
 \qquad |(\nabla\dd F_z)_x|\leq Cz.
\end{equation}
The last bound follows by differentiating the
smooth map $(y,v,z)\mapsto\exp_y(zv)$ twice in $y$, including the
first and second derivatives of $V$.  At $z=0$ the map is the
identity and its covariant second derivative vanishes.  Its
$z$-derivative is uniformly bounded for the bounded jets just
constructed, giving the displayed $Cz$ bound.  The exact middle
identity, rather than a near-isometry estimate, is essential.

The map $F_z$ takes $D$ locally into $D_z$, since its defining
geodesic from any $y\in D$ has length $z$.  Consequently
$\psi=\phi\circ F_z$ is an exterior support to $D$ at $x$.
Its derivative is $P^*\dd\phi$, so its gradient length is exactly
$|\nabla\phi(x')|$ and, by
Equation~\eqref{n3:domains:radial-normal}, its normal is
$\dot\gamma(0)$.  The covariant chain rule gives
\[
 \Hess\psi(X,Y)
   =\Hess\phi(PX,PY)
        +\dd\phi\bigl((\nabla\dd F_z)(X,Y)\bigr)
 \quad\hbox{at }x.
\]
Trace over an orthonormal basis of $\dot\gamma(0)^\perp$ and divide
by the common gradient length.  The first term is exactly the
corresponding normalized tangential trace at $x'$, and the remaining
term has absolute value at most $Cz$.  Finally, the tensor trace
changes by at most $Cz$, by the $C^1$ bound on $Q$ and parallel
transport.  It follows that
\[
 \bigl|\mathcal E_Q[\psi](x)
             -\mathcal E_Q[\phi](x')\bigr|\leq Cz.
\]
The assumed support inequality at $x$ proves
Equation~\eqref{n3:domains:parallel-estimate}.  There is no term
containing $z|\Hess\phi|$.
\end{proof}

\begin{lemma}[A smooth enclosure from weak supports]
\label{n3:domains:weak-support}
Let $N$ be a smooth compact connected three-manifold with boundary
$\Gamma_-\sqcup\Gamma_+$, where $\Gamma_+$ is nonempty and
$\Gamma_-$ may be empty.  Let $D$ be a nonempty closed subset which
contains a relative collar of $\Gamma_-$ when that boundary is
prescribed, and stays a positive distance from $\Gamma_+$.  Suppose
its interior frontier has exterior-support expansion at most $c$
for a smooth tensor $Q$.

For every $c'>c$ such that
$\Theta_Q(\Gamma_+)>c'$ with its normal out of $N$, there is a
smooth bounding domain $E$ containing $D$ in its relative interior,
containing the required inner collar, and avoiding $\Gamma_+$,
whose free boundary satisfies
\begin{equation}\label{n3:domains:smooth-enclosure-expansion}
 \Theta_Q(\partial E\setminus\Gamma_-)=c'.
\end{equation}
Its boundary and the domain may be disconnected.  In particular $D$
is contained in the closed total trapped region at any such larger
threshold.  No regularity or positive-thickness assumption on $D$
is required.
\end{lemma}

\begin{proof}
All distance neighborhoods in this proof are relative to $N$.
Since $D$ contains a collar of $\Gamma_-$ and avoids $\Gamma_+$,
their new frontiers, at sufficiently small radii, lie in a compact
subset of the interior.  A minimizing segment ending on a new
frontier has its nearest point on the interior frontier of $D$;
it cannot start on $\Gamma_-$ because its covered collar has fixed
positive width.  Thus the proof of
Lemma~\ref{n3:domains:parallel-support} applies unchanged.  Choose
$\beta>0$ so small that its conclusion holds for $0<z\leq\beta$,
all these frontiers avoid the prescribed boundary, and
\begin{equation}\label{n3:domains:band-gap}
 c+C\beta<c'.
\end{equation}

We first construct a smooth inner enclosure without imposing any
expansion bound on it.  Smoothly approximate the continuous distance
function $d_D=\operatorname{dist}(\cdot,D)$ on $N$ uniformly, with
error less than $\beta/32$, using a finite coordinate cover,
mollification, and a partition of unity.  Choose a regular value in
$(\beta/3,\beta/2)$ of the resulting smooth function.  Its sublevel
set $U$ has smooth interior boundary and, for example, satisfies
\begin{equation}\label{n3:domains:initial-smooth-neighborhood}
 D_{\beta/4}\subset U\subset D_{3\beta/4}.
\end{equation}
The boundary $\Gamma_-$ is contained in its relative interior and
$\Gamma_+$ is excluded.  If a component of $N\setminus\overline U$
meets no component of $\Gamma_+$, fill it.  This only deletes
components of the smooth free boundary.  After this filling the
second inclusion in
Equation~\eqref{n3:domains:initial-smooth-neighborhood} need not hold
for its interior, but every remaining free boundary component still
lies in the indicated thin neighborhood, and the first inclusion
still holds.

Choose $\eta\in C^\infty(N)$, $0\leq\eta\leq1$, equal to one
on that free boundary and supported in $\{d_D<\beta\}$.  Its
existence uses the positive separation between this compact smooth
boundary and the closed set $\{d_D\geq\beta\}$; it requires no
smoothness of $d_D$.  For a sufficiently large finite constant $A$,
the tensor
\begin{equation}\label{n3:domains:modified-tensor}
 Q_A=Q-\tfrac12(c'+A\eta)g
\end{equation}
makes the free boundary of $U$ a strictly negative inner barrier:
\[
 \Theta_{Q_A}(\partial U)
       =\Theta_Q(\partial U)-c'-A<0.
\]
All prescribed outer faces remain strictly positive for $Q_A$,
because $\eta=0$ there and their original expansion exceeds $c'$.
Apply the barrier part of
Theorem~\ref{n3:domains:region-theorem} on each component of
$N\setminus\overline U$.  Each meets $\Gamma_+$ by the filling
step and has nonempty inner boundary by connectedness of $N$ and
$U\ne\varnothing$.  We obtain smooth MOTSs $\Sigma$ bounding a
domain $E$ which contains $\overline U$ and avoids $\Gamma_+$.
They satisfy, with their normal out of $E$,
\begin{equation}\label{n3:domains:modified-mots-expansion}
 \Theta_Q(\Sigma)=c'+A\eta\geq c'.
\end{equation}

We claim $\Sigma$ misses $\{d_D<\beta\}$.  Otherwise set
$z=\min_{\Sigma}d_D<\beta$.  The first inclusion in
Equation~\eqref{n3:domains:initial-smooth-neighborhood} gives
$z\geq\beta/4>0$.  Moreover $D_z$ lies on the inner side of
$\Sigma$.  To see this directly, if $d_D(y)<z$, a shortest path
from $y$ to $D$ of length less than $z$ cannot cross $\Sigma$:
such a crossing would have distance from $D$ less than $z$.
Since $D\subset E$, the whole path and $y$ lie in $E$.
Closure gives $D_z\subset\overline E$.

At a point attaining the minimum, a defining function for the
smooth boundary $\Sigma$ is therefore an exterior support to
$D_z$, with the same outward normal as in
Equation~\eqref{n3:domains:modified-mots-expansion}.  The parallel
support bound and Equation~\eqref{n3:domains:band-gap} imply
\[
 \Theta_Q(\Sigma)\leq c+Cz<c',
\]
contradicting Equation~\eqref{n3:domains:modified-mots-expansion}.
This excludes the entire modification range, so $\eta=0$ on
$\Sigma$ and proves Equation~\eqref{n3:domains:smooth-enclosure-expansion}.
The contained $D_{\beta/4}$ gives strict enclosure of $D$.
If desired, fill any remaining complementary components which meet
no designated outer face; this deletes smooth boundary components
and preserves all stated properties.
\end{proof}

For an exterior, truncate by a distant sphere that is strict at the
larger threshold.  In the black construction, $\mathcal B$ already
contains the original inner collar.  In the white construction there is
no required inner boundary, and the black faces belong to $\Gamma_+$.
The later applications therefore need compact avoidance of the foreign
faces; they need no boundary regularity of the weak set.  Empty weak sets
can be omitted.

\subsection{Small offsets and the prepared domain}

\begin{proposition}[Prepared domain]\label{n3:domains:prepared-domain}
The thresholds and the domain above can be chosen with the following
properties.  The numbers $c_b,c_w<0$ are right-continuity points of
their respective full maximal-region families.  The boundary of
$\Omega$ consists of the faces in
Equation~\eqref{n3:domains:face-expansions}, has disjoint smooth
outward leaf collars as in
Lemma~\ref{n3:domains:stable-collar}, and preserves the enclosing-area
lower bound in Equation~\eqref{n3:domains:inherited-area}.

There are a smooth compactly supported function $C$ on $\Omega$,
constants $C_b\geq0$, $C_w\leq0$, and positive constants $k_B$
on its face collars such that
\begin{equation}\label{n3:domains:offset-collars}
 \begin{aligned}
  C_w\leq C\leq C_b,&\qquad
  C=C_b-k_Bs&&\text{in a black collar near }B_b,\\
  &\qquad C=C_w+k_Bs&&\text{in a white collar near }B_w.
 \end{aligned}
\end{equation}
For an absent type, its corresponding constant is simply a bound.
Set
\begin{equation}\label{n3:domains:assigned-heights}
 b_-=c_b-C_b<0,
 \qquad b_+=-c_w-C_w>0.
\end{equation}
There is a smooth positive function $\rho$, equal to a positive
constant times $r^{-4}$ on the far end and satisfying symbol
derivative bounds there, such that, for every $h\in[b_-,b_+]$,
\begin{equation}\label{n3:domains:offset-smallness}
 |(h+C)\tr K|+|\nabla C|+\rho
      \leq\tfrac12\mathfrak m
 \quad\hbox{everywhere on }\Omega.
\end{equation}
The thresholds, collars, offset, and $\rho$ are fixed before the
deformation parameters are chosen.  After fixing the thresholds and
collars, the offset and its first derivatives can be made arbitrarily
small.
\end{proposition}

\begin{proof}
The large-sphere comparison above confines all negative-threshold
regions to a fixed compact set, independently of how close the
thresholds are to zero.  Choose a slightly larger compact set
$\mathcal K$ containing this range and $\supp K$ in its interior.
Let
\[
 m_0=\min_{\mathcal K}\mathfrak m>0,
 \qquad T_0=\sup_{\mathcal K}|\tr K|.
\]
First restrict both negative thresholds to an interval so close to
zero that
\[
 \max\{|c_b|,|c_w|\}\,T_0<m_0/8,
\]
while retaining Equation~\eqref{n3:domains:black-interval}.  Lemma~\ref{n3:domains:right-continuity}
allows the required successive choices in that interval.  Construct
the two regions, select $\Omega$, and fix their disjoint collars
inside $\mathcal K$ as above.  All ensuing constants may depend on
these fixed choices.

Here is an explicit offset construction.  On each collar choose a
smooth nonincreasing cutoff $\chi(s)$, equal to one near $s=0$ and
zero near the collar's far end.  Choose $\kappa_B>0$ so that
$1-\kappa_Bs>0$ throughout that collar.  For a common amplitude
$a>0$, define the contribution of a black collar to be
$a\chi(s)(1-\kappa_Bs)$, and that of a white collar to be
$-a\chi(s)(1-\kappa_Bs)$.  Extend by zero outside the disjoint
collars and sum.  This produces a smooth function on the manifold
with boundary $\Omega$, supported in $\mathcal K$, with
$C_b=a$, $C_w=-a$, and $k_B=a\kappa_B$ on the respective collars.
These remain valid bounds when a type is absent.  Its $C^1$ norm
tends to zero with $a$, even though the collar widths have already
been fixed.

For $h\in[b_-,b_+]$ and $C\in[C_w,C_b]$,
\[
 |h+C|\leq\max\{|c_b|,|c_w|\}+C_b-C_w.
\]
Choose $a$ so small that
$(C_b-C_w)T_0+\|\nabla C\|_\infty<m_0/8$.
The first two terms of
Equation~\eqref{n3:domains:offset-smallness} are then at most
$m_0/4$ on $\mathcal K$ and vanish outside $\mathcal K$.
Finally take a fixed smooth positive weight equal to $r^{-4}$ far
out and multiply it by a sufficiently small positive constant.
The end lower bound for $\mathfrak m$ and $\delta<1$, together with
compact positivity, allow the choice $\rho\leq\mathfrak m/4$
everywhere.  This proves Equation~\eqref{n3:domains:offset-smallness}.
All remaining assertions were proved in constructing the regions
and collars.
\end{proof}

At the assigned face heights the quantity $h+C$ consequently has
the useful exact values
\begin{equation}\label{n3:domains:compatible-face-values}
 \begin{aligned}
  h+C&=c_b=P_B+H&&\text{on }B_b\text{ when }h=b_-,\\
  h+C&=-c_w=P_B-H&&\text{on }B_w\text{ when }h=b_+.
 \end{aligned}
\end{equation}
These identities determine the deformation's boundary signs.
Right-continuity and Lemma~\ref{n3:domains:weak-support} will exclude
limiting sublevel sets away from the corresponding full trapped regions.
